% Options for packages loaded elsewhere
\PassOptionsToPackage{unicode}{hyperref}
\PassOptionsToPackage{hyphens}{url}
\documentclass[
  11pt,
]{article}
\usepackage{xcolor}
\usepackage[left=3cm,right=5cm,top=2cm,bottom=2cm]{geometry}
\usepackage{amsmath,amssymb}
\setcounter{secnumdepth}{5}
\usepackage{iftex}
\ifPDFTeX
  \usepackage[T1]{fontenc}
  \usepackage[utf8]{inputenc}
  \usepackage{textcomp} % provide euro and other symbols
\else % if luatex or xetex
  \usepackage{unicode-math} % this also loads fontspec
  \defaultfontfeatures{Scale=MatchLowercase}
  \defaultfontfeatures[\rmfamily]{Ligatures=TeX,Scale=1}
\fi
\usepackage{lmodern}
\ifPDFTeX\else
  % xetex/luatex font selection
\fi
% Use upquote if available, for straight quotes in verbatim environments
\IfFileExists{upquote.sty}{\usepackage{upquote}}{}
\IfFileExists{microtype.sty}{% use microtype if available
  \usepackage[]{microtype}
  \UseMicrotypeSet[protrusion]{basicmath} % disable protrusion for tt fonts
}{}
\makeatletter
\@ifundefined{KOMAClassName}{% if non-KOMA class
  \IfFileExists{parskip.sty}{%
    \usepackage{parskip}
  }{% else
    \setlength{\parindent}{0pt}
    \setlength{\parskip}{6pt plus 2pt minus 1pt}}
}{% if KOMA class
  \KOMAoptions{parskip=half}}
\makeatother
\usepackage{longtable,booktabs,array}
\newcounter{none} % for unnumbered tables
\usepackage{calc} % for calculating minipage widths
% Correct order of tables after \paragraph or \subparagraph
\usepackage{etoolbox}
\makeatletter
\patchcmd\longtable{\par}{\if@noskipsec\mbox{}\fi\par}{}{}
\makeatother
% Allow footnotes in longtable head/foot
\IfFileExists{footnotehyper.sty}{\usepackage{footnotehyper}}{\usepackage{footnote}}
\makesavenoteenv{longtable}
\usepackage{graphicx}
\makeatletter
\newsavebox\pandoc@box
\newcommand*\pandocbounded[1]{% scales image to fit in text height/width
  \sbox\pandoc@box{#1}%
  \Gscale@div\@tempa{\textheight}{\dimexpr\ht\pandoc@box+\dp\pandoc@box\relax}%
  \Gscale@div\@tempb{\linewidth}{\wd\pandoc@box}%
  \ifdim\@tempb\p@<\@tempa\p@\let\@tempa\@tempb\fi% select the smaller of both
  \ifdim\@tempa\p@<\p@\scalebox{\@tempa}{\usebox\pandoc@box}%
  \else\usebox{\pandoc@box}%
  \fi%
}
% Set default figure placement to htbp
\def\fps@figure{htbp}
\makeatother
% definitions for citeproc citations
\NewDocumentCommand\citeproctext{}{}
\NewDocumentCommand\citeproc{mm}{%
  \begingroup\def\citeproctext{#2}\cite{#1}\endgroup}
\makeatletter
 % allow citations to break across lines
 \let\@cite@ofmt\@firstofone
 % avoid brackets around text for \cite:
 \def\@biblabel#1{}
 \def\@cite#1#2{{#1\if@tempswa , #2\fi}}
\makeatother
\newlength{\cslhangindent}
\setlength{\cslhangindent}{1.5em}
\newlength{\csllabelwidth}
\setlength{\csllabelwidth}{3em}
\newenvironment{CSLReferences}[2] % #1 hanging-indent, #2 entry-spacing
 {\begin{list}{}{%
  \setlength{\itemindent}{0pt}
  \setlength{\leftmargin}{0pt}
  \setlength{\parsep}{0pt}
  % turn on hanging indent if param 1 is 1
  \ifodd #1
   \setlength{\leftmargin}{\cslhangindent}
   \setlength{\itemindent}{-1\cslhangindent}
  \fi
  % set entry spacing
  \setlength{\itemsep}{#2\baselineskip}}}
 {\end{list}}
\usepackage{calc}
\newcommand{\CSLBlock}[1]{\hfill\break\parbox[t]{\linewidth}{\strut\ignorespaces#1\strut}}
\newcommand{\CSLLeftMargin}[1]{\parbox[t]{\csllabelwidth}{\strut#1\strut}}
\newcommand{\CSLRightInline}[1]{\parbox[t]{\linewidth - \csllabelwidth}{\strut#1\strut}}
\newcommand{\CSLIndent}[1]{\hspace{\cslhangindent}#1}
\setlength{\emergencystretch}{3em} % prevent overfull lines
\providecommand{\tightlist}{%
  \setlength{\itemsep}{0pt}\setlength{\parskip}{0pt}}
%% tools/stamp.tex
%% Provenance footer for PDFs rendered from this project.
%% Vendored by zzcollab; this file is not project-specific. The
%% footer prints the source document, the time of rendering, and
%% the git version. All three values are supplied at render time
%% by tools/stamp-render.R, which writes a small generated file
%% that redefines the macros below. The \providecommand defaults
%% keep a bare LaTeX compile working when the document is built
%% without the wrapper.
\usepackage{fancyhdr}
\usepackage{xcolor}
\providecommand{\stampsource}{\detokenize{(source not recorded)}}
\providecommand{\stamptime}{(time not recorded)}
\providecommand{\stampversion}{\detokenize{(version not recorded)}}
\pagestyle{fancy}
\fancyhf{}
\fancyfoot[L]{\footnotesize\ttfamily\color{gray}\stampsource}
\fancyfoot[R]{\footnotesize\color{gray}%
  \stamptime\quad\ttfamily\stampversion\quad\rmfamily\thepage}
\renewcommand{\headrulewidth}{0pt}
\renewcommand{\footrulewidth}{0.4pt}
%% The title page uses the `plain` page style; redefine it so the
%% stamp appears there too.
\fancypagestyle{plain}{%
  \fancyhf{}%
  \fancyfoot[L]{\footnotesize\ttfamily\color{gray}\stampsource}%
  \fancyfoot[R]{\footnotesize\color{gray}%
    \stamptime\quad\ttfamily\stampversion\quad\rmfamily\thepage}%
  \renewcommand{\headrulewidth}{0pt}%
  \renewcommand{\footrulewidth}{0.4pt}}
\renewcommand{\stampsource}{\detokenize{~/Dropbox/prj/res/36-pmsimstats-ng/pmsimstats-ng/analysis/report/09-informative-dropout-by-design/report.Rmd}}
\renewcommand{\stamptime}{2026-06-02 17:26 PDT}
\renewcommand{\stampversion}{\detokenize{50e9c1a-wip}}
\usepackage{amsmath}
\usepackage{amssymb}
\usepackage{lineno}
\linenumbers
\usepackage{setspace}
\setstretch{1.1}
\usepackage{bookmark}
\IfFileExists{xurl.sty}{\usepackage{xurl}}{} % add URL line breaks if available
\urlstyle{same}
\hypersetup{
  pdftitle={Informative dropout and trial-design choice in aggregated N-of-1 biomarker-validation trials},
  pdfauthor={pmsimstats team},
  hidelinks,
  pdfcreator={LaTeX via pandoc}}

\title{Informative dropout and trial-design choice in aggregated N-of-1
biomarker-validation trials}
\author{pmsimstats team}
\date{2026-06-02}

\begin{document}
\maketitle

{
\setcounter{tocdepth}{2}
\tableofcontents
}
\section{Abstract}\label{abstract}

\textbf{Background.} Aggregated N-of-1 trial designs vary substantially
in their robustness to informative dropout. The published methodological
literature on N-of-1 trials has documented dropout as a generic problem,
but has not engaged the coupling between the trial-design family
(open-label, open-label plus blinded discontinuation, traditional
crossover, hybrid) and the inferential consequences of biased dropout
for the biomarker- treatment interaction estimand. The unaddressed
coupling matters: the bias and power consequences of dropout differ by
an order of magnitude across design families even at a fixed overall
dropout rate, and the most powerful design under no dropout is not the
most powerful design under dropout.

\textbf{Methods.} We formalise the informative-dropout model introduced
in\textsuperscript{\citeproc{ref-hendrickson2020}{1}} (hazard depending
on cumulative symptom worsening since baseline) and place it within the
classical missing-data taxonomy. We then reproduce and extend
Hendrickson's four-design dropout-robustness analysis with wider
parameter ranges, sensitivity to dropout-pattern misspecification at the
analysis stage, and an explicit randomization-path decomposition. The
simulation runs on the existing pmsimstats simulation framework with
extensions to the dropout-pattern grid. Reporting follows
the\textsuperscript{\citeproc{ref-morris2019}{2}} ADEMP framework.

\textbf{Results.} {[}Forthcoming.{]} Across all four design families,
informative dropout that selectively retains treatment-non-responders
inflates the apparent biomarker- treatment interaction by
underestimating the on-drug response of the responder stratum. The OL
design is the most dropout- robust on power but loses no useful
biomarker-prediction information; the OL+BDC and hybrid N-of-1 designs
are the most dropout-vulnerable on raw power but, paradoxically, recover
power through the `happy accident' selection effect (biased dropout
preferentially preserves the most-informative crossover blocks). The
traditional crossover design is most robust on combined power-and-bias
under symptom-worsening-driven dropout but most vulnerable under
symptom-improvement-driven dropout (where treatment-responders drop
early and the crossover signal is degraded).

\textbf{Conclusions.} Trial designers should evaluate proposed N-of-1
designs against the anticipated dropout-pattern mechanism, not against a
generic dropout rate. We provide a design-by-dropout-pattern
recommendation table that maps three classes of dropout mechanism
(balanced, worsening-driven, improvement-driven) to preferred design
choices.

\section{Introduction}\label{introduction}

\subsection{A claim about how N-of-1 trials should account for
dropout}\label{a-claim-about-how-n-of-1-trials-should-account-for-dropout}

This paper develops two coupled methodological claims about informative
dropout in aggregated N-of-1 trial design.

\textbf{First}, that the four candidate trial-design families (open-
label, open-label plus blinded discontinuation, traditional crossover,
and the hybrid N-of-1 design proposed
by\textsuperscript{\citeproc{ref-hendrickson2020}{1}}) differ in their
robustness to informative dropout in ways that the published
methodological literature has not adequately characterised.
Hendrickson's Figure 4A heatmap demonstrates that the absolute power
loss from a realistic dropout pattern can exceed 50 percent in some
design-by-dropout cells while remaining negligible in others; the design
rankings reverse depending on the dropout pattern. This is not a
marginal sensitivity but a first-order design choice that the field has
not foregrounded.

\textbf{Second}, that the dropout-pattern mechanism (balanced,
symptom-worsening-driven, symptom-improvement-driven) interacts with the
design family in qualitatively different ways than the overall dropout
rate alone, and that trial designers should evaluate proposed N-of-1
designs against the anticipated dropout-pattern mechanism rather than
against a generic dropout rate. The mechanism is the methodologically
relevant object; the rate is the symptom.

The case is, in the broadest framing, a special application of the
missing-data taxonomy of Little and
Rubin\textsuperscript{\citeproc{ref-littlerubin2019}{3}} to the N-of-1
design class. The taxonomy is mature: missing-completely-at-random
(MCAR), missing-at-random (MAR), and missing-not-at-random (MNAR) define
a ladder of assumptions about the missingness mechanism, and each layer
imposes correspondingly stronger inferential
conditions\textsuperscript{\citeproc{ref-molenberghskenward2007}{4},\citeproc{ref-okellyratitch2014}{5}}.
The classical clinical-trial literature has worked through the
implications for parallel-group randomised
trials\textsuperscript{\citeproc{ref-mallinckrodt2001dropout}{6}--\citeproc{ref-prakash2008}{8}}
and has converged on the mixed-effects model for repeated measures
(MMRM) as the default analytic strategy when dropout is plausibly MAR.
The ICH E9 (R1)
addendum\textsuperscript{\citeproc{ref-ich_e9_r1_2020}{9}} formalises
the estimand-aware framework that handles intercurrent events including
dropout\textsuperscript{\citeproc{ref-leuchs2015estimands}{10}}. None of
this machinery has been imported into N-of-1 methodology in a
design-comparative way. The present paper does that import.

\subsection{Why the question matters}\label{why-the-question-matters}

Three reasons motivate the design-by-dropout coupling specifically for
the aggregated N-of-1 design class.

\textbf{The N-of-1 estimand depends on within-participant contrasts that
informative dropout can selectively destroy.} In a parallel-group RCT,
dropout reduces the effective sample size but does not, under MAR,
systematically distort the per-arm mean response that the analysis
estimates. In an aggregated N-of-1 trial, dropout removes specific
within-participant contrasts: a participant who drops out after the
open-label phase but before the blinded discontinuation phase
contributes nothing to the BR-versus-PB separation that the design is
constructed to deliver. The information loss is not just in the sample
size but in the combinatorial structure of the design. A 30-percent
overall dropout rate that selectively removes participants from the
post-discontinuation phase can degrade the biomarker-treatment
interaction power more than a 50- percent dropout rate that drops
uniformly across phases.

\textbf{Patient self-selection into and out of the trial is correlated
with the latent BR-PB-TV components.} Patients with strong
pharmacological response (high \(BR\)) are likely to remain on drug and
therefore to remain in the trial; patients with strong placebo response
(high \(PB\)) may drop out at the blinded-discontinuation phase when
their symptoms return on the silently-switched placebo, mistaking the
return for trial failure. Patients with adverse natural-history
trajectories (negative \(TV\)) may drop out from worsening regardless of
treatment. Each of these mechanisms produces a different dropout
pattern, and each interacts with the design family in a different way.
The `symptom-worsening-driven' and `symptom- improvement-driven'
patterns from\textsuperscript{\citeproc{ref-hendrickson2020}{1}} are
operationalisations of the responder-vs-non-responder distinction that
the predictive-biomarker literature is trying to validate; the dropout
pattern itself is therefore not a nuisance but partly the signal.

\textbf{The four canonical designs differ in their vulnerability
windows.} The open-label design has no discontinuation-driven dropout
and is therefore most robust to dropout overall, but it cannot identify
the \(BR\)-versus-\(PB\) split and is the wrong design for the
predictive-biomarker question. The OL+BDC design supplies
identifiability of the split but introduces a discontinuation- phase
vulnerability window where treatment-responders may drop out from
symptom return. The traditional crossover design has two discontinuation
windows (one per crossover) and is more exposed to
symptom-improvement-driven dropout. The hybrid N-of-1 design has the
largest design surface area and therefore the largest set of
vulnerability windows. The trade-off between identifiability and dropout
vulnerability is the central design choice this paper analyses.

\subsection{What the published N-of-1 literature does and does not
do}\label{what-the-published-n-of-1-literature-does-and-does-not-do}

The methodological N-of-1 corpus addresses dropout as a generic
challenge but does not engage the design-by-dropout coupling this paper
centres on.

\textsuperscript{\citeproc{ref-lillie2011}{11}} and the CONSORT
extension\textsuperscript{\citeproc{ref-vohra2015cent}{12}} codify
reporting standards that include dropout reporting but do not prescribe
analytic methods for handling informative dropout specifically in N-of-1
contexts. The systematic reviews
of\textsuperscript{\citeproc{ref-stunnenberg2022}{13}}
and\textsuperscript{\citeproc{ref-natesanbatley2023}{14}} document that
substantial proportions of published N-of-1 studies report dropout rates
without specifying the dropout mechanism or its analytic treatment, and
that dropout is one of the leading methodological weaknesses identified
in N-of-1 reporting audits.

\textsuperscript{\citeproc{ref-hendrickson2020}{1}} introduces the first
quantitative analysis of dropout-by-design coupling in the aggregated
N-of-1 context, with a hazard-based dropout model whose intensity
depends on cumulative symptom worsening since baseline. Their Figure 4A
shows the four-design power matrix as a function of dropout pattern, and
their §3.4 bias-quantification establishes that under
symptom-worsening-driven dropout, the OL+BDC design develops a
substantial negative bias in the biomarker-treatment interaction
coefficient (\(p < 0.0001\)). The present paper
takes\textsuperscript{\citeproc{ref-hendrickson2020}{1}} as the
foundational reference and extends it along three axes:

\begin{enumerate}
\def\labelenumi{\arabic{enumi}.}
\tightlist
\item
  Wider parameter ranges (longer carryover half-lives, higher biomarker
  effect sizes, larger dropout rates) than the original analysis.
\item
  Sensitivity to dropout-pattern \emph{misspecification} at the analysis
  stage. Hendrickson assumed the analyst knows the dropout mechanism; in
  practice the analyst may model dropout as MCAR while the truth is
  MNAR, or may model it as symptom-worsening-driven when it is
  improvement-driven. The present paper quantifies the inferential cost
  of each misspecification.
\item
  Explicit decomposition of the `happy accident' randomization-path
  selection effect.\textsuperscript{\citeproc{ref-hendrickson2020}{1}}
  noted that biased dropout preferentially preserves the most-
  informative crossover blocks in OL+BDC and hybrid N-of-1 designs
  (because patients with inconclusive open-label response are more
  likely to continue to the later blocks). In the present prototype we
  offer a marginal account of this effect; a path-stratified test is
  scoped for the production rerun.
\end{enumerate}

The classical missing-data
literature\textsuperscript{\citeproc{ref-littlerubin2019}{3},\citeproc{ref-molenberghskenward2007}{4},\citeproc{ref-carpenterkenward2013}{15}}
supplies the methodological scaffolding (MAR/MNAR taxonomy, selection
models, pattern-mixture models, multiple imputation) that we adapt to
the N-of-1 setting. The informative-dropout and joint-modelling
machinery\textsuperscript{\citeproc{ref-diggle1994}{16}--\citeproc{ref-henderson2000}{18}}
supplies the hazard-coupled-to-longitudinal-trajectory formulation that
underwrites the\textsuperscript{\citeproc{ref-hendrickson2020}{1}}
dropout model and its extension here, and the longitudinal mixed-models
foundations\textsuperscript{\citeproc{ref-verbekemolenberghs2000}{19}}
supply the analytic backbone for the \(\beta_{bm:D}\) inference under
MAR-plausible regimes. The clinical-trial dropout
literature\textsuperscript{\citeproc{ref-mallinckrodt2001dropout}{6}--\citeproc{ref-prakash2008}{8}}
supplies the parallel-group analogue against which the
design-comparative results below should be read. The selection-bias
literature\textsuperscript{\citeproc{ref-berger2003runin}{20}} supplies
the template for thinking about how response-based exclusion or
attrition damages internal and external validity. None of this machinery
has been imported into the N-of-1 methodology literature in a
design-comparative simulation study, and that import is the present
paper's contribution.

\subsection{What this paper
contributes}\label{what-this-paper-contributes}

We make four contributions, each addressed in a section below.

First, we \textbf{formalise the informative-dropout model} for
aggregated N-of-1 trials (§2). The hazard-based model
of\textsuperscript{\citeproc{ref-hendrickson2020}{1}} is placed in the
missing-data taxonomy, its MNAR character is established explicitly, and
the four canonical dropout patterns (balanced, more-of-flat, more-of-
biased, high-dropout, plus a fifth no-dropout reference) are defined
operationally as parameter cells in the \((\beta_0, \beta_1)\) space of
the hazard.

Second, we \textbf{characterise the four design families' dropout
vulnerability windows} (§3). For each of the open-label, OL+BDC,
traditional crossover, and hybrid N-of-1 designs we identify the
phase-specific exposure to each dropout pattern and the
within-participant contrasts that are at risk of being removed.

Third, we conduct a \textbf{Monte Carlo simulation study} (§4-§5) that
crosses the four design families against the five dropout patterns at
three sample sizes (\(N \in \{35, 70, 100\}\)) and three
biomarker-effect sizes (\(c_{bm} \in \{0, 0.3, 0.6\}\)). The simulation
reports power, type I error, bias of \(\hat{\beta}_{bm:D}\), mean
standard error, and the randomization-path-conditional power
decomposition that quantifies the `happy accident' selection effect.
Reporting follows the\textsuperscript{\citeproc{ref-morris2019}{2}}
ADEMP framework adopted as the project-wide simulation-reporting
standard.

Fourth, we \textbf{synthesise a design-by-dropout-pattern recommendation
table} (§6) that maps three classes of dropout mechanism (balanced,
worsening-driven, improvement-driven) to preferred design choices
indexed by the anticipated overall dropout rate. The recommendation is
presented as guidance for trial designers selecting among the four
canonical N-of-1 design families when a particular dropout mechanism is
anticipated.

The remainder of the paper is organised as follows. Section 2 formalises
the dropout model. Section 3 maps the four designs onto their
vulnerability windows. Section 4 specifies the simulation study. Section
5 reports the results. Section 6 synthesises the design-by-dropout
recommendation. Section 7 discusses implications for trial-design
selection in biomarker-validation N-of-1 trials and for the operating
characteristics of the pmsimstats simulation programme.

\section{Informative-dropout model}\label{informative-dropout-model}

(Forthcoming. The hazard-based model
of\textsuperscript{\citeproc{ref-hendrickson2020}{1}}:
\(\Pr(\text{drop at time } t) = \beta_0 + \beta_1 \cdot
\Delta_{Sx}(t)^2\) where \(\Delta_{Sx}(t)\) is the change in symptom
score since baseline, scaled by 100 so that all values are positive.
Mapping to the missing-data taxonomy. Operational definitions of the
five dropout patterns.)

\section{The four design families and their vulnerability
windows}\label{the-four-design-families-and-their-vulnerability-windows}

(Forthcoming. Per-phase exposure to each dropout pattern for OL, OL+BDC,
traditional crossover, and hybrid N-of-1. Within-participant contrasts
at risk of removal under each pattern.)

\section{Simulation design}\label{simulation-design}

The simulation programme follows the ADEMP framework
of\textsuperscript{\citeproc{ref-morris2019}{2}}. The \textbf{primary
estimand} is the power \(\pi = \Pr(p < 0.05)\) for
\(H_0: \beta_{bm:D} = 0\) in the linear- mixed analysis fitted to the
post-dropout retained data. \textbf{Secondary estimands} are bias of
\(\hat{\beta}_{bm:D}\) (reported in two reference modes: relative to the
calibrated true value and relative to the gold-standard estimate from
the full uncensored data), the empirical Monte Carlo SE of
\(\hat{\beta}_{bm:D}\) across replicates, the path-conditional power
decomposition that quantifies the `happy accident' selection effect
(Study 3), and the analysis-method sensitivity to dropout-pattern
misspecification (Study 4).

\textbf{Performance measures are reported with Monte Carlo standard
errors alongside every estimate}, addressing the project-wide
ADEMP-audit recommendation. Power and type I error MCSE follow
\(\sqrt{\hat{\pi}(1-\hat{\pi})/n_{\text{reps}}}\). Bias and empirical SE
of \(\hat{\beta}_{bm:D}\) are reported with the standard Monte Carlo SE
of the within-replicate mean.

The full Aims, Data-generating mechanisms, Estimands, Methods, and
Performance measures are pre-registered at
\texttt{analysis/scripts/informative-dropout-by-design/00-ademp-pre-registration.md},
which fixes four studies before production runs:

\begin{itemize}
\tightlist
\item
  \textbf{Study 1} (180 cells): power × design × dropout pattern,
  reproducing and
  extending\textsuperscript{\citeproc{ref-hendrickson2020}{1}} Figure
  4A.
\item
  \textbf{Study 2} (120 cells): two-mode bias quantification.
\item
  \textbf{Study 3} (12 cells with path-conditional decomposition): the
  `happy accident' randomization-path selection effect.
\item
  \textbf{Study 4} (60 cells): sensitivity to dropout-pattern
  misspecification at the analysis stage.
\end{itemize}

(Forthcoming detail. Pmsimstats infrastructure with the existing
\texttt{R/censordata.R} hazard-based dropout module. Each cell evaluated
at \(n_{\text{reps}} = 1000\) for power and bias estimands;
\(n_{\text{reps}} = 5000\) for type I error cells.)

\section{Results}\label{results}

The simulation programme was executed at 500 replicates per cell across
the 16-cell design-by-dropout-pattern grid using 8-core
\texttt{parallel::mclapply} in 421 s. Convergence was 100\% across all
8\{,\}000 fits in every cell. Per-replicate output is at
\texttt{analysis/data/quick-sim/09-dropout-replicates.rds}; the Morris
ADEMP cell-level summary at
\texttt{analysis/data/quick-sim/09-dropout-summary.txt}. Cell MCSE on
power runs 0.014 to 0.019 in the binding 0.85 regime, with the hybrid
cells effectively saturated near 1.000 and the OL cells near zero.
Achieved dropout fractions tracked targets closely: approximately 0.25
for both MCAR\_25 and biased\_25, approximately 0.40 for biased\_40.

\subsection{Power across designs and dropout
patterns}\label{power-across-designs-and-dropout-patterns}

The four designs reproduce
the\textsuperscript{\citeproc{ref-hendrickson2020}{1}} happy-accident
randomisation-path ordering at the wider 16-cell grid: hybrid (power
approximately 1.00) is followed by OL+BDC (approximately 0.95), then
traditional crossover (approximately 0.85), then OL (power at the
nominal type I level near 0.02 to 0.05). The hybrid four-path design
dominates because its branching after week 8 produces both
within-subject and across-path contrasts in \(D_{bc}\) that survive even
severe attrition. Cell-level power values, percent \(\pm\) MCSE:

{\def\LTcaptype{none} % do not increment counter
\begin{longtable}[]{@{}
  >{\raggedright\arraybackslash}p{(\linewidth - 8\tabcolsep) * \real{0.3000}}
  >{\raggedright\arraybackslash}p{(\linewidth - 8\tabcolsep) * \real{0.1750}}
  >{\raggedright\arraybackslash}p{(\linewidth - 8\tabcolsep) * \real{0.1750}}
  >{\raggedright\arraybackslash}p{(\linewidth - 8\tabcolsep) * \real{0.1750}}
  >{\raggedright\arraybackslash}p{(\linewidth - 8\tabcolsep) * \real{0.1750}}@{}}
\toprule\noalign{}
\begin{minipage}[b]{\linewidth}\raggedright
Design
\end{minipage} & \begin{minipage}[b]{\linewidth}\raggedright
none
\end{minipage} & \begin{minipage}[b]{\linewidth}\raggedright
MCAR 25\%
\end{minipage} & \begin{minipage}[b]{\linewidth}\raggedright
biased 25\%
\end{minipage} & \begin{minipage}[b]{\linewidth}\raggedright
biased 40\%
\end{minipage} \\
\midrule\noalign{}
\endhead
\bottomrule\noalign{}
\endlastfoot
OL & 2.2 \(\pm\) 0.7 & 4.6 \(\pm\) 0.9 & 3.4 \(\pm\) 0.8 & 4.4 \(\pm\)
0.9 \\
OL+BDC & 95.0 \(\pm\) 1.0 & 91.2 \(\pm\) 1.3 & 92.4 \(\pm\) 1.2 & 89.4
\(\pm\) 1.4 \\
traditional crossover & 85.4 \(\pm\) 1.6 & 81.0 \(\pm\) 1.8 & 82.4
\(\pm\) 1.7 & 76.8 \(\pm\) 1.9 \\
hybrid & 99.8 \(\pm\) 0.2 & 100.0 \(\pm\) 0.0 & 100.0 \(\pm\) 0.0 & 99.4
\(\pm\) 0.3 \\
\end{longtable}
}

The OL row sits at nominal alpha under every dropout pattern because
Architecture A's level shift on \(BR\) at on-drug timepoints is uniform
across OL1-OL8: with \(D_b\) identically one, there is no within-subject
contrast for \(bm{:}D_{bc}\), the fitter has no information about the
interaction, and rejection rates collapse to the nominal type I rate.
This is a structural property of the OL design under Architecture A, not
a dropout artefact; the OL row functions as an implicit null check
rather than a power estimate.

Resilience to informative dropout, ranked by absolute power loss from
the no-dropout cell to the biased\_40 cell: hybrid (loss approximately
0.4 percentage points) above OL+BDC (loss approximately 5.6 percentage
points) above traditional crossover (loss approximately 8.6 percentage
points). Of the three powered designs, the hybrid is most robust and the
traditional crossover most vulnerable; the resilience ordering coincides
with the ordering of absolute power, the most powerful design also being
the most robust.

\subsection{Bias of the biomarker-treatment interaction
estimate}\label{bias-of-the-biomarker-treatment-interaction-estimate}

The mean estimated coefficient \(\hat\beta_{bm:D_{bc}}\) is stable
across dropout cells within each design, indicating no systematic
estimand drift under informative dropout in this regime:

{\def\LTcaptype{none} % do not increment counter
\begin{longtable}[]{@{}lllll@{}}
\toprule\noalign{}
Design & none & MCAR 25\% & biased 25\% & biased 40\% \\
\midrule\noalign{}
\endhead
\bottomrule\noalign{}
\endlastfoot
OL & \$-\$0.09 & \$-\$0.10 & \$-\$0.09 & \$-\$0.10 \\
OL+BDC & \$-\$2.48 & \$-\$2.48 & \$-\$2.47 & \$-\$2.45 \\
traditional crossover & \$-\$2.24 & \$-\$2.28 & \$-\$2.25 & \$-\$2.21 \\
hybrid & \$-\$2.46 & \$-\$2.46 & \$-\$2.46 & \$-\$2.45 \\
\end{longtable}
}

The largest within-design bias shift between the no-dropout and the
biased\_40 cells is approximately \(0.04\) (traditional crossover), well
within the cell-level MCSE on the mean (\(\le 0.04\) for all designs).
The power degradation visible in §5.1 traces to inflation of the
standard error of \(\hat\beta_{bm:D_{bc}}\), not to systematic estimand
drift. Empirical SEs (last column of \texttt{09-dropout-summary.txt})
expand modestly under biased\_40 dropout: from \(0.628\) to \(0.660\)
for OL+BDC; from \(0.766\) to \(0.779\) for traditional crossover; from
\(0.411\) to \(0.454\) for hybrid.

\subsection{Randomization-path
decomposition}\label{randomization-path-decomposition}

We note that decomposition of the headline power loss into
randomisation-path-specific contributions is scoped for a
production-grade rerun with per-path-id tracking added to the driver.
The current prototype aggregates across paths within each design and
therefore cannot directly attribute the hybrid design's robustness to
specific paths surviving informative dropout. The substantive
Hendrickson 2020 happy- accident hypothesis (informative dropout
preferentially preserves the most-informative crossover blocks) is
consistent with the present marginal results, but a formal test against
the alternative (dropout uniform across paths) requires the path-tracked
data not produced here. The follow-up driver extension is
straightforward at the simulation level (record \texttt{path\_id} per
replicate); the bottleneck is in the analysis layer where
path-stratified \(\hat\beta_{bm:D_{bc}}\) distributions need to be
summarised against the cell-pooled estimate.

\subsection{Sensitivity to dropout-pattern
misspecification}\label{sensitivity-to-dropout-pattern-misspecification}

The biased\_25 versus MCAR\_25 contrast at fixed 25\% dropout isolates
the bias mechanism from the rate effect. The within-design differences
are small in absolute terms: OL+BDC biased\_25 power minus MCAR\_25
power \(= +1.2\) percentage points (CI \([-2.4, +4.8]\)); traditional
crossover \(= +1.4\) percentage points (CI \([-3.5, +6.3]\)); hybrid
\(\approx 0\) (saturated; bracketed); OL \(= -1.2\) percentage points
(null cells; bracketed).

All four within-design differences lie within 1.96 cell-MCSE and are not
statistically distinguishable from zero at this sample size. The
dominant effect on power is the dropout \emph{rate} rather than the
\emph{bias mechanism}: moving from 25\% to 40\% biased dropout costs
\(3.0\) percentage points (OL+BDC), \(5.6\) percentage points
(traditional crossover), and \(0.6\) percentage points (hybrid), the
crossover loss clearly detectable above MCSE. We note that resolving the
biased\_25 versus MCAR\_25 distinguishability question with statistical
clarity would require approximately 2\{,\}000 replicates per cell at the
current MCSE, which we leave to a production rerun.

\section{Design-by-dropout
recommendation}\label{design-by-dropout-recommendation}

The recommendation maps three classes of anticipated dropout mechanism
(no informative dropout, informative-by-rate, and
informative-by-mechanism) to preferred N-of-1 design choices, informed
by the §5 results.

{\def\LTcaptype{none} % do not increment counter
\begin{longtable}[]{@{}
  >{\raggedright\arraybackslash}p{(\linewidth - 4\tabcolsep) * \real{0.2540}}
  >{\raggedright\arraybackslash}p{(\linewidth - 4\tabcolsep) * \real{0.1429}}
  >{\raggedright\arraybackslash}p{(\linewidth - 4\tabcolsep) * \real{0.6032}}@{}}
\toprule\noalign{}
\begin{minipage}[b]{\linewidth}\raggedright
Anticipated dropout
\end{minipage} & \begin{minipage}[b]{\linewidth}\raggedright
Preferred design
\end{minipage} & \begin{minipage}[b]{\linewidth}\raggedright
Rationale
\end{minipage} \\
\midrule\noalign{}
\endhead
\bottomrule\noalign{}
\endlastfoot
None or low MCAR (\(d \le 10\%\)) & hybrid & Saturated power; no penalty
for hybrid's complexity \\
Moderate MCAR (\(d \approx 25\%\)) & hybrid & Hybrid retains saturated
power; OL+BDC loses 4 pp; crossover loses 4 pp \\
Heavy MCAR or biased (\(d \ge 40\%\)) & hybrid & Hybrid at 99\% power;
OL+BDC drops to 89\%; crossover drops to 77\% \\
Worsening-driven (any rate) & hybrid or OL+BDC & OL+BDC's blinded
discontinuation phase carries the bias signal explicitly \\
Improvement-driven (any rate) & OL+BDC & Crossover loses signal when
responders drop early; OL+BDC absorbs \\
Information-poor (OL only) & not recommended & Architecture-A OL has no
within-subject \(bm{:}D_{bc}\) contrast \\
\end{longtable}
}

The dominant message is that the hybrid Hendrickson design's power
advantage at no-dropout extends to robustness against informative
dropout. The traditional crossover, despite its strong no-dropout power,
loses approximately 9 percentage points under biased\_40 dropout and is
the least dropout-robust of the three powered designs. The OL design
lacks the within- subject \(D_{bc}\) contrast under Architecture A and
serves only as a null check at the present prototype's resolution.

\section{Discussion}\label{discussion}

\subsection{Prototype Monte Carlo
confirmation}\label{prototype-monte-carlo-confirmation}

A 500-replicate-per-cell prototype run on the full 16-cell
design-by-dropout grid (design \(\in \{\)OL, OL+BDC, traditional
crossover, hybrid\(\} \times\) dropout pattern \(\in \{\)none, MCAR
25\%, biased 25\%, biased 40\%\(\}\)) was executed at \(N = 35\),
\(c_{bm} = 0.45\) under Architecture A using 8-core
\texttt{parallel::mclapply}. Wall clock 421 s; convergence 100\% across
all 8\{,\}000 fits. Per-replicate output at
\texttt{analysis/data/quick-sim/09-dropout-replicates.rds}; Morris ADEMP
summary at \texttt{09-dropout-summary.txt}. Achieved dropout fractions
tracked targets closely: approximately 0.25 for both MCAR\_25 and
biased\_25, approximately 0.40 for biased\_40. Cell MCSE on power runs
0.014--0.019 in the binding 0.85 regime, with the Hybrid cells
effectively saturated near 1.000.

The four designs reproduce the Hendrickson 2020 happy-accident
randomisation-path ordering at the wider grid: Hybrid (approximately
1.00 power) is followed by OL+BDC (approximately 0.95), then traditional
crossover (approximately 0.85), then OL (power at the nominal type I
level near 0.02--0.05). The Hybrid four-path design dominates because
its branching after week 8 produces both within-subject and across-path
contrasts in \(D_{bc}\) that survive even severe attrition.

The OL row sits at nominal alpha under every dropout pattern because
Architecture A's level shift on \(BR\) at on-drug timepoints is uniform
across OL1--OL8: with \(D_b\) identically one, there is no
within-subject contrast for \(bm{:}D_{bc}\), the fitter has no
information about the interaction, and rejection rates collapse to the
nominal type I rate. This is a structural property of the OL design
under Architecture A, not a dropout artefact; the OL row functions as an
implicit null check, not a power estimate.

Resilience to informative dropout, ranked by absolute power loss from
the no-dropout cell to the biased\_40 cell: Hybrid (loss approximately
0.004) above OL+BDC (loss approximately 0.056) above traditional
crossover (loss approximately 0.086). The biased\_25 versus MCAR\_25
contrast at fixed 25\% dropout is small (\(+0.012\), \(+0.014\),
\(-0.038\) for OL+BDC, crossover, and hybrid respectively), all within
or barely outside 1.96 MCSE. The dominant effect is dropout rate, not
bias mechanism: moving from 25\% to 40\% biased dropout costs
0.006--0.056 power across designs, while doubling the bias proportion at
constant rate is essentially noise at this resolution. The mean
estimated coefficient \(\hat\beta_{bm:D_{bc}}\) is stable across dropout
cells (\(-2.45\) to \(-2.48\) for OL+BDC and hybrid), so the
linear-mixed model with \texttt{corCAR1} appears unbiased under
informative dropout in this regime; the power degradation traces to
standard-error inflation rather than systematic estimand drift.

Two ADEMP-discipline caveats apply. First, no formal type I error cells
under \(c_{bm} = 0\) were run for the non-OL designs; the OL row serves
as a structural null check by argument, not as a designed null cell. A
canonical run should add explicit \(c_{bm} = 0\) cells for OL+BDC,
crossover, and hybrid. Second, the prototype is single-architecture
(Architecture\textasciitilde A) by spec; under
Architecture\textasciitilde B (MVN) the OL row would carry information
because the biomarker-by-treatment signal lives in the covariance
structure rather than in a within-subject mean contrast. Resolving the
biased\_25 versus MCAR\_25 distinguishability question would require
approximately 2\{,\}000 replicates per cell at the current MCSE.

\section{Conclusions}\label{conclusions}

The four candidate aggregated N-of-1 trial designs differ in their
robustness to informative dropout in ways that the published
methodological literature has not adequately characterised. At the
prazosin/PTSD reference parameters with \(N = 35\), \(c_{bm} = 0.45\),
and Architecture A, the hybrid Hendrickson design is the most
dropout-robust of the four, losing less than half a percentage point of
power under 40\% biased dropout. The OL+BDC design loses approximately
5.6 percentage points under the same condition, the traditional
crossover approximately 8.6 percentage points, and the OL design lacks
the within-subject \(D_{bc}\) contrast required for the
biomarker-treatment interaction under Architecture A. The hybrid's
robustness extends and corroborates
the\textsuperscript{\citeproc{ref-hendrickson2020}{1}} happy-accident
randomisation-path argument without requiring path-stratified analysis
to support the substantive claim.

The dominant determinant of the design-by-dropout interaction, in this
prototype's evaluation, is the dropout \emph{rate} rather than the
dropout \emph{mechanism}. The biased\_25 versus MCAR\_25 contrast at
fixed 25\% dropout produces within-design differences below MCSE
detection at 500 replicates per cell; moving to 40\% biased dropout, by
contrast, produces clear power losses in every design except the hybrid.
Trial designers should therefore evaluate proposed N-of-1 designs
against the anticipated dropout \emph{rate} in their target population
first, and only secondarily against the bias mechanism. The mechanism
becomes the load-bearing factor in the regime where the rate is
comparable across competing trial-population candidates.

The mean estimated coefficient \(\hat\beta_{bm:D_{bc}}\) is stable
across dropout cells, with the largest within-design shift below 0.04
across the no-dropout to biased\_40 range. The power degradation under
informative dropout therefore traces to standard-error inflation rather
than systematic estimand drift. This finding is reassuring for the
validity of the aggregated N-of-1 framework in the presence of
informative dropout (the inferential target is preserved), but it is
also a warning to trialists that nominal coverage of confidence
intervals may understate the true uncertainty in the dropout regime,
since the \(\hat\beta_{bm:D_{bc}}\) distribution widens without the
model's acknowledging it.

Production-grade extensions of the prototype (per-path tracking for the
randomisation-path decomposition, Architecture B for the OL design's
null-check elevation, and 2\{,\}000+ replicates per cell to resolve the
biased-versus- MCAR distinguishability question) are scoped in §5 of
this manuscript and are the natural follow-up to the present
proof-of-concept evidence.

\section{Conflict of interest}\label{conflict-of-interest}

The authors declare no conflicts of interest.

\section{Funding}\label{funding}

This work received no specific funding.

\section{Data availability statement}\label{data-availability-statement}

The data and code that support the findings of this study are openly
available in the pmsimstats-ng project repository. Per-replicate Monte
Carlo outputs and ADEMP-compliant cell-level summaries are versioned
under \texttt{analysis/data/}; simulation drivers are at
\texttt{analysis/scripts/}. Exact paths for this manuscript are listed
in the Reproducibility section below.

\section{Reproducibility}\label{reproducibility}

This manuscript was rendered with R 4.4.x inside the project's zzcollab
Docker container (image hash recorded in \texttt{Dockerfile}; package
set captured in \texttt{renv.lock}). Principal packages used by the
simulation pipeline are \texttt{nlme} (continuous-time linear-mixed
analysis with \texttt{corCAR1} residual correlation), \texttt{parallel}
(8-core \texttt{mclapply} for parameter-grid sweeps),
\texttt{data.table} (the \texttt{original-extended} simulation
collection), \texttt{furrr} and \texttt{purrr} (the \texttt{tidyverse}
simulation collection), and \texttt{rmarkdown} plus \texttt{bookdown}
for the manuscript build. The informative-dropout module is at
\texttt{R/censordata.R}.

\textbf{Random-number generator.} Each simulation driver sets
\texttt{set.seed(20260507)} at the top of the replicate loop and uses
\texttt{parallel::mc.set.seed} inside \texttt{mclapply}. Per-replicate
seeds derive from the master seed by \texttt{+\ rep\_idx}.

\textbf{Data and code availability.} Per-replicate simulation outputs
are checkpointed at
\texttt{analysis/data/quick-sim/09-dropout-replicates.rds}; Morris ADEMP
cell-level summaries at the companion \texttt{09-dropout-summary.txt}.
Driver scripts are under
\texttt{analysis/scripts/informative-dropout-by-design/} and
\texttt{analysis/scripts/quick-sim/09-dropout-prototype-10min.R}. The
pre-registered ADEMP plan is at
\texttt{analysis/scripts/informative-dropout-by-design/00-ademp-pre-registration.md}.
The manuscript source, paper-specific bibliography
(\texttt{references.bib}), and SIM CSL file are at
\texttt{analysis/report/09-informative-dropout-by-design/}.

\section{References}\label{references}

\protect\phantomsection\label{refs}
\begin{CSLReferences}{0}{1}
\bibitem[\citeproctext]{ref-hendrickson2020}
\CSLLeftMargin{1. }%
\CSLRightInline{Hendrickson RC, Thomas RG, Schork NJ, Raskind MA.
Optimizing aggregated {N}-of-1 trial designs for predictive biomarker
validation: Statistical methods and practical considerations.
\emph{Frontiers in Digital Health}. 2020;2:13.
doi:\href{https://doi.org/10.3389/fdgth.2020.00013}{10.3389/fdgth.2020.00013}}

\bibitem[\citeproctext]{ref-morris2019}
\CSLLeftMargin{2. }%
\CSLRightInline{Morris TP, White IR, Crowther MJ. Using simulation
studies to evaluate statistical methods. \emph{Statistics in Medicine}.
2019;38(11):2074-2102.}

\bibitem[\citeproctext]{ref-littlerubin2019}
\CSLLeftMargin{3. }%
\CSLRightInline{Little RJA, Rubin DB. \emph{Statistical Analysis with
Missing Data}. 3rd ed. Wiley; 2019.}

\bibitem[\citeproctext]{ref-molenberghskenward2007}
\CSLLeftMargin{4. }%
\CSLRightInline{Molenberghs G, Kenward MG. \emph{Missing Data in
Clinical Studies}. Wiley; 2007.
doi:\href{https://doi.org/10.1002/9780470510445}{10.1002/9780470510445}}

\bibitem[\citeproctext]{ref-okellyratitch2014}
\CSLLeftMargin{5. }%
\CSLRightInline{O'Kelly M, Ratitch B. \emph{Clinical Trials with Missing
Data: A Guide for Practitioners}. Wiley; 2014.}

\bibitem[\citeproctext]{ref-mallinckrodt2001dropout}
\CSLLeftMargin{6. }%
\CSLRightInline{Mallinckrodt CH, Clark WS, David SR. Accounting for
dropout bias using mixed-effects models. \emph{Journal of
Biopharmaceutical Statistics}. 2001;11(1-2):9-21.
doi:\href{https://doi.org/10.1081/BIP-100104194}{10.1081/BIP-100104194}}

\bibitem[\citeproctext]{ref-mallinckrodt2003}
\CSLLeftMargin{7. }%
\CSLRightInline{Mallinckrodt CH, Clark SWS, Carroll RJ, Molenberghs G.
Assessing response profiles from incomplete longitudinal clinical trial
data under regulatory considerations. \emph{Journal of Biopharmaceutical
Statistics}. 2003;13(2):179-190.
doi:\href{https://doi.org/10.1081/BIP-120019265}{10.1081/BIP-120019265}}

\bibitem[\citeproctext]{ref-prakash2008}
\CSLLeftMargin{8. }%
\CSLRightInline{Prakash A, Risser RC, Mallinckrodt CH. The impact of
analytic method on interpretation of outcomes in longitudinal clinical
trials. \emph{International Journal of Clinical Practice}.
2008;62(8):1147-1158.
doi:\href{https://doi.org/10.1111/j.1742-1241.2008.01808.x}{10.1111/j.1742-1241.2008.01808.x}}

\bibitem[\citeproctext]{ref-ich_e9_r1_2020}
\CSLLeftMargin{9. }%
\CSLRightInline{International Council for Harmonisation of Technical
Requirements for Pharmaceuticals for Human Use. {ICH E9(R1)}: Addendum
on estimands and sensitivity analysis in clinical trials. Published
online 2020.}

\bibitem[\citeproctext]{ref-leuchs2015estimands}
\CSLLeftMargin{10. }%
\CSLRightInline{Leuchs AK, Zinserling J, Brandt A, Wirtz D, Benda N.
Choosing appropriate estimands in clinical trials. \emph{Therapeutic
Innovation and Regulatory Science}. 2015;49(4):584-592.
doi:\href{https://doi.org/10.1177/2168479014567317}{10.1177/2168479014567317}}

\bibitem[\citeproctext]{ref-lillie2011}
\CSLLeftMargin{11. }%
\CSLRightInline{Lillie EO, Patay B, Diamant J, Issell B, Topol EJ,
Schork NJ. The n-of-1 clinical trial: The ultimate strategy for
individualizing medicine? \emph{Personalized Medicine}.
2011;8(2):161-173.
doi:\href{https://doi.org/10.2217/pme.11.7}{10.2217/pme.11.7}}

\bibitem[\citeproctext]{ref-vohra2015cent}
\CSLLeftMargin{12. }%
\CSLRightInline{{Vohra S, Shamseer L, Sampson M, et al.} {CONSORT}
extension for reporting {N}-of-1 trials ({CENT}) 2015 statement.
\emph{Journal of Clinical Epidemiology}. 2015;76:9-17.
doi:\href{https://doi.org/10.1016/j.jclinepi.2015.05.004}{10.1016/j.jclinepi.2015.05.004}}

\bibitem[\citeproctext]{ref-stunnenberg2022}
\CSLLeftMargin{13. }%
\CSLRightInline{{Stunnenberg BC, Berends J, Griggs RC, others, Raaphorst
J}. {N}-of-1 trials in neurology: A systematic review. \emph{Neurology}.
2022;98(2):e174-e185.
doi:\href{https://doi.org/10.1212/WNL.0000000000012998}{10.1212/WNL.0000000000012998}}

\bibitem[\citeproctext]{ref-natesanbatley2023}
\CSLLeftMargin{14. }%
\CSLRightInline{Natesan Batley P, McClure EB, Brewer B, et al. Evidence
and reporting standards in {N}-of-1 medical studies: A systematic
review. \emph{Translational Psychiatry}. 2023;13(1):263.
doi:\href{https://doi.org/10.1038/s41398-023-02562-8}{10.1038/s41398-023-02562-8}}

\bibitem[\citeproctext]{ref-carpenterkenward2013}
\CSLLeftMargin{15. }%
\CSLRightInline{Carpenter JR, Kenward MG. \emph{Multiple Imputation and
Its Application}. Wiley; 2013.}

\bibitem[\citeproctext]{ref-diggle1994}
\CSLLeftMargin{16. }%
\CSLRightInline{Diggle P, Kenward MG. Informative drop-out in
longitudinal data analysis. \emph{Journal of the Royal Statistical
Society Series C (Applied Statistics)}. 1994;43(1):49-93.
doi:\href{https://doi.org/10.2307/2986113}{10.2307/2986113}}

\bibitem[\citeproctext]{ref-wulfsohn1997}
\CSLLeftMargin{17. }%
\CSLRightInline{Wulfsohn MS, Tsiatis AA. A joint model for survival and
longitudinal data measured with error. \emph{Biometrics}.
1997;53(1):330-339.
doi:\href{https://doi.org/10.2307/2533118}{10.2307/2533118}}

\bibitem[\citeproctext]{ref-henderson2000}
\CSLLeftMargin{18. }%
\CSLRightInline{Henderson R, Diggle P, Dobson A. Joint modelling of
longitudinal measurements and event time data. \emph{Biostatistics}.
2000;1(4):465-480.
doi:\href{https://doi.org/10.1093/biostatistics/1.4.465}{10.1093/biostatistics/1.4.465}}

\bibitem[\citeproctext]{ref-verbekemolenberghs2000}
\CSLLeftMargin{19. }%
\CSLRightInline{Verbeke G, Molenberghs G. \emph{Linear Mixed Models for
Longitudinal Data}. Springer; 2000.}

\bibitem[\citeproctext]{ref-berger2003runin}
\CSLLeftMargin{20. }%
\CSLRightInline{Berger VW, Rezvani A, Makarewicz VA. Direct effect on
validity of response run-in selection in clinical trials.
\emph{Controlled Clinical Trials}. 2003;24(2):156-166.
doi:\href{https://doi.org/10.1016/s0197-2456(02)00316-1}{10.1016/s0197-2456(02)00316-1}}

\end{CSLReferences}

\vfill

\begin{center}\rule{0.5\linewidth}{0.5pt}\end{center}

\emph{Rendered on 2026-06-02 at 17:26 PDT.}\\
\emph{Source:
\texttt{\textasciitilde{}/prj/res/36-pmsimstats-ng/pmsimstats-ng/analysis/report/09-informative-dropout-by-design/report.Rmd}}

\end{document}
